\section{Présentation du corpus d'étude}

\subsection{Choix du corpus}
Cette étude porte sur 31 romans d'Émile Zola, publiés entre 1865 et 1903. Les nouvelles, les essais et les textes journalistiques ont été écartés afin de constituer un ensemble relativement homogène du point de vue du genre et de la longueur. Les romans offrent également un contexte suffisamment développé pour étudier les thèmes, les milieux sociaux et les motifs récurrents de l'œuvre zolienne.

Ce choix répond aussi à une contrainte méthodologique. L'intégration de textes courts aurait accentué les différences de volume entre les œuvres et aurait pu déséquilibrer l'analyse thématique. La restriction du corpus au genre romanesque ne supprime pas toutes les variations de longueur, mais elle les limite et facilite la comparaison entre les textes.

Le corpus réunit cinq romans antérieurs aux \textit{Rougon-Macquart}, les vingt volumes de ce cycle, les trois romans des \textit{Trois Villes} et les trois volumes publiés des \textit{Quatre Évangiles}. Dans ce dernier cycle, \textit{Vérité} paraît à titre posthume, tandis que \textit{Justice}, resté à l'état d'ébauche, n'a jamais été publié.

\begin{longtable}{r p{8cm} r}
\caption{Nombre de tokens par roman de Zola dans l’ordre chronologique}
\label{tab:tokens_romans_zola} \\

\toprule
\textbf{Date} & \textbf{Titre du roman} & \textbf{Nombre de tokens} \\
\midrule
\endfirsthead

\toprule
\textbf{Date} & \textbf{Titre du roman} & \textbf{Nombre de tokens} \\
\midrule
\endhead

1865 & \textit{La confession de Claude} & 47253 \\
1866 & \textit{Le vœu d'une morte} & 38722 \\
1867 & \textit{Les mystères de Marseille} & 125253 \\
1867 & \textit{Thérèse Raquin} & 66363 \\
1868 & \textit{Madeleine Férat} & 97318 \\
1871 & \textit{La fortune des Rougon} & 116685 \\
1872 & \textit{La curée} & 105120 \\
1873 & \textit{Le ventre de Paris} & 110442 \\
1874 & \textit{La conquête de Plassans} & 113245 \\
1875 & \textit{La faute de l'abbé Mouret} & 114667 \\
1876 & \textit{Son Excellence Eugène Rougon} & 126186 \\
1877 & \textit{L'assommoir} & 158340 \\
1878 & \textit{Une page d'amour} & 102185 \\
1880 & \textit{Nana} & 141573 \\
1882 & \textit{Pot-Bouille} & 134844 \\
1883 & \textit{Au Bonheur des dames} & 147521 \\
1884 & \textit{La joie de vivre} & 117255 \\
1885 & \textit{Germinal} & 164473 \\
1886 & \textit{L'œuvre} & 130613 \\
1887 & \textit{La terre} & 163687 \\
1888 & \textit{Le rêve} & 66379 \\
1890 & \textit{La bête humaine} & 125418 \\
1891 & \textit{L'argent} & 143010 \\
1892 & \textit{La débâcle} & 187046 \\
1893 & \textit{Le docteur Pascal} & 112713 \\
1894 & \textit{Lourdes} & 173511 \\
1896 & \textit{Rome} & 233958 \\
1898 & \textit{Paris} & 177389 \\
1899 & \textit{Fécondité} & 220131 \\
1901 & \textit{Travail} & 198712 \\
1903 & \textit{Vérité} & 223274 \\

\bottomrule
\end{longtable}

Ce corpus couvre ainsi l'ensemble de la carrière romanesque de Zola, de ses premiers textes à ses dernières œuvres. Cette amplitude chronologique permet d'observer les régularités thématiques de son écriture, mais aussi leurs transformations éventuelles au fil du temps, qui seront étudiées dans le chapitre consacré à BERTopic.

\subsection{Préparation des textes}

Les romans ont été réunis au format \texttt{.txt}, puis nettoyés afin de retirer les éléments éditoriaux qui ne relèvent pas directement du récit : mentions d'édition, tables des matières, notes, numéros de page et intitulés de chapitres. Cette étape vise à éviter que ces éléments n'influencent artificiellement les résultats\footnote{Pour ce qui est des scripts de nettoyage, ils sont disponible en annexe et/ou sur mon github}.

Après nettoyage, chaque fichier a été organisé de manière à ne conserver qu'une phrase par ligne. Cette organisation ne modifie pas le contenu des romans, mais facilite leur traitement automatique et leur segmentation. Elle permet également d'appliquer la même procédure à des textes provenant initialement de formats différents, notamment \texttt{TXT}, \texttt{PDF} ou \texttt{EPUB}.

Le nombre de tokens varie sensiblement d'un roman à l'autre. Les premiers textes, comme \textit{Le vœu d'une morte}, sont relativement courts, tandis que des œuvres tardives comme \textit{Rome}, \textit{Fécondité}, \textit{Travail} ou \textit{Vérité} sont beaucoup plus volumineuses. Cette différence doit être prise en compte, car un roman long produit mécaniquement davantage d'occurrences lexicales et d'unités d'analyse qu'un roman court.

\subsection{Préparation du corpus pour l'analyse thématique}

Les romans ont été découpés en segments de 40 phrases. Cette taille constitue un compromis entre deux objectifs : conserver suffisamment de contexte pour identifier un thème et produire des unités plus fines que les chapitres ou les œuvres entières.

Plusieurs tailles de segments ont été expérimentées, notamment 10, 40, 50, 70 et 100 phrases. Les segments de 40 phrases ont donné les résultats les plus lisibles et les plus cohérents au cours des essais exploratoires. Ce choix reste néanmoins en partie empirique et constitue donc une limite méthodologique.

Un découpage par chapitre ou par paragraphe a également été envisagé. Ces solutions ont été écartées en raison de la forte variation de longueur de ces unités. De plus, la structure des paragraphes n'est pas toujours conservée de manière fiable lors de la conversion des fichiers. Les dialogues produisent notamment de nombreux paragraphes très courts. Le découpage par phrases offre donc une solution plus homogène et plus facilement applicable à l'ensemble du corpus.

\subsection{Statistiques descriptives du corpus}

\begin{longtable}{p{8cm}r}
\caption{Nombre de paquets de phrases par roman} \\
\toprule
\textbf{Titre du roman} & \textbf{Nombre de paquets} \\
\midrule
\endfirsthead

\toprule
\textbf{Titre du roman} & \textbf{Nombre de paquets} \\
\midrule
\endhead

\textit{La confession de Claude} & 65 \\
\textit{Le vœu d'une morte} & 60 \\
\textit{Les mystères de Marseille} & 184 \\
\textit{Thérèse Raquin} & 80 \\
\textit{Madeleine Férat} & 124 \\
\textit{La fortune des Rougon} & 154 \\
\textit{La curée} & 123 \\
\textit{Le ventre de Paris} & 136 \\
\textit{La conquête de Plassans} & 165 \\
\textit{La faute de l'abbé Mouret} & 171 \\
\textit{Son Excellence Eugène Rougon} & 191 \\
\textit{L'assommoir} & 228 \\
\textit{Une page d'amour} & 166 \\
\textit{Nana} & 225 \\
\textit{Pot-Bouille} & 207 \\
\textit{Au Bonheur des dames} & 198 \\
\textit{La joie de vivre} & 164 \\
\textit{Germinal} & 225 \\
\textit{L'œuvre} & 163 \\
\textit{La terre} & 222 \\
\textit{Le rêve} & 82 \\
\textit{La bête humaine} & 154 \\
\textit{L'argent} & 155 \\
\textit{La débâcle} & 221 \\
\textit{Le docteur Pascal} & 135 \\
\textit{Lourdes} & 198 \\
\textit{Rome} & 231 \\
\textit{Paris} & 206 \\
\textit{Fécondité} & 259 \\
\textit{Travail} & 208 \\
\textit{Vérité} & 227 \\

\bottomrule
\end{longtable}

La segmentation produit un nombre variable de paquets selon la longueur des œuvres. Les romans les plus courts, comme \textit{Le vœu d'une morte} et \textit{La confession de Claude}, fournissent moins d'unités d'analyse, tandis que \textit{Fécondité}, \textit{Rome} ou \textit{L'assommoir} en produisent davantage.

\begin{table}[H]
\centering
\caption{Statistiques descriptives du nombre de paquets par roman}
\begin{tabular}{lr}
\toprule
\textbf{Indicateur} & \textbf{Valeur} \\
\midrule
Nombre de romans & 31 \\
Moyenne & 171.84 \\
Écart-type & 52.40 \\
Minimum & 60 \\
Premier quartile & 145 \\
Médiane & 171 \\
Troisième quartile & 214.5 \\
Maximum & 259 \\
\bottomrule
\end{tabular}
\end{table}

Le corpus comprend en moyenne environ 172 paquets par roman, avec une médiane de 171. Ces deux valeurs proches montrent que la distribution n'est pas fortement déséquilibrée. L'écart-type de 52 paquets et l'écart entre le minimum de 60 et le maximum de 259 témoignent toutefois d'une variation réelle entre les œuvres.

Cette différence doit être gardée à l'esprit lors de l'interprétation des résultats : les romans les plus longs contribuent à un plus grand nombre de segments et peuvent donc occuper une place plus importante dans les résultats globaux. La segmentation améliore la comparabilité des unités textuelles, mais elle ne supprime pas entièrement l'effet lié à la longueur des œuvres.


\begin{table}[H]
\centering
\caption{Statistiques descriptives du nombre de mots par paquet}
\begin{tabular}{lr}
\toprule
\textbf{Indicateur} & \textbf{Valeur} \\
\midrule
Nombre de paquets & 5327 \\
Moyenne & 785.30 \\
Écart-type & 212.13 \\
Minimum & 42 \\
Premier quartile & 636 \\
Médiane & 753 \\
Troisième quartile & 902.5 \\
Maximum & 2944 \\
\bottomrule
\end{tabular}
\end{table}

Au total, le corpus segmenté comprend 5,327 paquets. Ils contiennent en moyenne 785 mots, pour une médiane de 753 mots. La moitié centrale des segments se situe entre 636 et environ 903 mots, ce qui indique que la majorité des unités possède une taille relativement comparable.

La variation restante s'explique par le choix d'un nombre fixe de phrases plutôt que d'un nombre fixe de mots. Selon le rythme narratif et la longueur des phrases, deux segments peuvent donc présenter des volumes différents. Les valeurs extrêmes, comprises entre 42 et 2,944 mots, correspondent probablement à des segments résiduels ou à des passages composés de phrases particulièrement longues. Ces cas doivent être pris en compte, car les segments les plus volumineux contiennent davantage d'information lexicale.


\subsection{Lemmatisation du corpus de Zola}

Les segments ont ensuite été lemmatisés\footnote{Le code est disponible sur mon gitbub} avec la bibliothèque Python \texttt{spaCy} et son modèle français \texttt{fr\_core\_news\_lg}. La lemmatisation rassemble les différentes formes fléchies d'un mot sous une forme commune : « mange », « manges » et « manger » sont ainsi associés au lemme « manger ». Cette normalisation réduit la dispersion du vocabulaire et facilite la construction de la matrice terme-document.

Une liste d'exclusion personnalisée a été constituée au fil des essais. Elle comprend les mots-outils, la ponctuation, les nombres et certains termes très fréquents mais peu utiles à l'interprétation des thèmes. Plusieurs noms de personnages ont également été retirés, car leur présence tendait à produire des topics centrés sur les protagonistes plutôt que sur des ensembles thématiques plus généraux.

Pour chaque segment, les lemmes sont convertis en minuscules, puis filtrés. Seuls les noms et les adjectifs sont conservés, ces catégories étant particulièrement utiles pour faire apparaître des contenus thématiques tels que le travail, la famille, la religion, l'argent ou les rapports sociaux. Ce choix oriente donc volontairement l'analyse vers les thèmes et non vers le style grammatical de Zola.

La liste d'exclusion a été ajustée après l'examen des termes les plus fréquents et des premiers résultats. Cette procédure empirique améliore la lisibilité des topics, mais elle introduit également une part d'intervention humaine qui doit être prise en compte dans leur interprétation. À l'issue du traitement, chaque roman est représenté par un ensemble de segments lemmatisés et filtrés, utilisés pour les analyses par NMF et K-means.
